\clearpage\fancyhead[L]{\small PALLAS.}\hypertarget{native-354-lower}{}\label{native-354-lower}
\begin{center}\textbf{PALLAS.}\par\vspace{4pt}\hrule\vspace{4pt}\fontsize{10}{12}\selectfont Göttingische gelehrte Anzeigen. Heft 67, S. 657\;..\;660, 25. April 1812.\par\vspace{4pt}\hrule\vspace{4pt}\end{center}

\par\indent In diesen Blättern haben wir schon einige Male von der Arbeit des Professors \textsc{Gauß} über die Störungen der Pallas berichtet (siehe die Hefte~8 und 130 des Vorjahres [S.~322 und 327 dieses Bandes]). Die dort angedeuteten Ergebnisse dieser Untersuchungen haben nun erneut eine schöne Bestätigung gefunden. Am 4. April erlaubte das Wetter zum ersten Mal in diesem Jahr, den Planeten aufzusuchen. Genau an der Stelle, die die von Herrn \textsc{Nicolai} nach der Störungstheorie berechnete Ephemeride angab, war ein Sternchen elfter Größe sichtbar. Am 9. April hatte es diese Stelle verlassen und sich wie erwartet weiterbewegt. Ein paar Vergleiche mit 93 Herculis ergaben die folgende Stellung, wobei die Deklination allerdings nur geschätzt ist.

\begin{center}\fontsize{10}{12}\selectfont\setlength{\tabcolsep}{0pt}\renewcommand{\arraystretch}{1.15}\begin{tabular}{l}
April 9. $11^{\mathrm u}\;32^{\mathrm m}\;27^{\mathrm s}$ M. Z.\\
Rektaszension der Pallas $269^\circ\;6'\;6''7$\\
Nördliche Deklination $16\;38$
\end{tabular}\end{center}
\SourceLeaf{355}\Needspace{13\baselineskip}\par\indent Demnach stimmt die Stellung des Planeten vollkommen mit jener Ephemeride überein. Gleich darauf erfuhr Professor \textsc{Gauß} von Herrn Dr. \textsc{Olbers}, den er gebeten hatte, den Planeten aufzusuchen, dass ihm das schon am 3. April gelungen war. Die beiden folgenden Bremer Beobachtungen stimmen ebenfalls aufs Genaueste mit der Ephemeride von Herrn \textsc{Nicolai} überein.

\begin{center}
\fontsize{10}{12}\selectfont\setlength{\tabcolsep}{3pt}\renewcommand{\arraystretch}{1.15}\begin{tabular}{c|c|c}
Mittlere Zeit & \shortstack{Rektaszension} & \shortstack{Nördliche\\Deklination} \\
\hline
April 3 \quad $11^{\mathrm u}\;26^{\mathrm m}\;44^{\mathrm s}$ & $268^\circ\;42'\;23''1$ & $15^\circ\;25'\;4''9$ \\
\phantom{April} 4 \qquad\quad $12\;2\;33$ & $268\;47\;12.5$ & $15\;37\;41.1$ \\
\end{tabular}
\end{center}

\par\indent Wir beeilen uns, diese Stellungen bekannt zu machen. Der Planet ist wegen seiner Lichtschwäche in diesem Jahr äußerst schwer zu beobachten. Wir bitten alle Astronomen dringend, ihn besonders um die Zeit der Opposition (10. Juni) sorgfältig zu beobachten.

\par\indent Auch die \emph{erste vorläufige} Berechnung der \emph{allgemeinen} Theorie der Pallasstörungen hat Professor \textsc{Gauß} vor Kurzem vollendet. Die Zahl aller periodischen Gleichungen, die aus der Einwirkung Jupiters folgen und noch merklich sind, erreicht \emph{Vier hundert}. Unter den sehr bemerkenswerten Ergebnissen, die sich schon daraus ableiten lassen, ist eines von ganz besonderem Interesse. Aus bestimmten Gründen halten wir es hier in der folgenden Chiffre fest. Den Schlüssel dazu werden wir zu gegebener Zeit mitteilen.

\begin{center}
1111000100101001
\end{center}

\par\indent Bei dieser Gelegenheit fügen wir noch zwei weitere interessante astronomische Mitteilungen aus einem Brief von Herrn Professor \textsc{Bessel} in Königsberg an Professor \textsc{Gauß} vom 26. März bei.

\par\indent Bekanntlich hat schon \textsc{Herschel} die Auffassung vertreten, dass die Sterne, die uns als Doppelsterne erscheinen, nicht zufällig mit uns auf einer geraden Linie liegen, sondern sich tatsächlich nahe beieinander befinden und jedes Paar ein eigenes System bildet. Gegen \emph{die Gründe}, aus denen dieser berühmte Beobachter den Schluss zog, bestehen zwar noch erhebliche Einwände, doch andere Beobachtungen sprechen sehr stark für die Auffassung selbst. Aus Gründen der Wahrscheinlichkeitsrechnung ist es nämlich höchst unwahrscheinlich, dass \emph{der Zufall} unter den Sternen bis zu einer gewissen Ordnung auch nur ein einziges Doppelsternpaar gebildet haben sollte, geschweige denn so viele, wie sich tatsächlich am Himmel finden. Nun hat diese Auffassung durch eine sehr interessante Bemerkung von Herrn Professor \textsc{Bessel} eine sehr einleuchtende Bestätigung erhalten. Bei seiner wichtigen Bearbeitung der Beobachtungen \textsc{Bradley}s, mit der er sich schon seit mehreren Jahren beschäftigt, hat er mehrere Fixsterne mit sehr starker Eigenbewegung erkannt. Am bemerkenswertesten darunter ist der Doppelstern 61 Cygni, dessen Eigenbewegung die stärkste aller bekannten Eigenbewegungen ist. Beide Sterne haben diese Bewegung gemeinsam. Sie ist auf den Herkules gerichtet. Auch alle starken Eigenbewegungen folgen dieser Richtung sehr genau. Für 1755 findet Herr \textsc{Bessel} aus \textsc{Bradley}s Beobachtungen die folgenden Werte.

\begin{center}
\fontsize{10}{12}\selectfont\setlength{\tabcolsep}{3pt}\renewcommand{\arraystretch}{1.15}\begin{tabular}{c|c|c}
 & \shortstack{Rektaszension} & \shortstack{Nördliche\\Deklination} \\
\hline
61 Cygni & $313^\circ\;59'\;11''3$ & $37^\circ\;33'\;29''7$ \\
Sequens & $313\;59\;25.7$ & $47\;33\;45.7$ \\
\end{tabular}
\end{center}

\par\indent Die Angabe \textsc{Piazzi}s konnte nicht verwendet werden, weil das Beobachtungsjahr unbekannt ist. Zwei Beobachtungen in der \emph{Histoire céleste} liefern jedoch die jährliche Bewegung des Hauptsterns

\begin{center}\fontsize{10}{12}\selectfont\setlength{\tabcolsep}{0pt}\renewcommand{\arraystretch}{1.15}\begin{tabular}{l}
$+5''250$ in Rektaszension;\\$+2''321$ in Deklination.
\end{tabular}\end{center}
\SourceLeaf{356}\noindent Könnte man die jährliche Parallaxe dieses Doppelsterns durch Beobachtungen bestimmen, so ließe sich künftig, sobald die Bewegung der beiden Sterne um ihren gemeinsamen Schwerpunkt bekannt ist, daraus das Verhältnis der Summe ihrer Massen zur Masse unserer Sonne ableiten. Daran sieht man, wie sehr dieser bemerkenswerte Stern die größte Aufmerksamkeit der Astronomen verdient.

\Needspace{25\baselineskip}\par\indent Eine zweite Mitteilung von Herrn Professor \textsc{Bessel} betrifft den großen Kometen des Vorjahres. Schon vor acht Monaten (siehe unsere Gelehrten Anzeigen vom Vorjahr, S.~1293 [S.~329 dieses Bandes]) bemerkten wir, dass Hoffnung bestehe, diesen Kometen im April und in den folgenden Monaten im Wassermann ein drittes Mal wiederzufinden. Es gibt noch keinen Grund, diese Hoffnung aufzugeben. Sofern seine Helligkeit sich nicht durch physische Ursachen ändert, wird der Komet noch ein Fünftel des Lichts haben, das er Anfang Januar hatte, als man ihn noch mit bloßem Auge erkannte. Die folgende kleine Ephemeride, die Herr Professor \textsc{Bessel} berechnet hat, wird bei der Suche gute Dienste leisten.

\begin{center}
\fontsize{10}{12}\selectfont\setlength{\tabcolsep}{3pt}\renewcommand{\arraystretch}{1.15}\begin{tabular}{c|c|c|c}
\shortstack{Mittlere Pariser Zeit\\$1812\;16^{\mathrm u}$} & \shortstack{Rektasz.} & \shortstack{Südliche\\Dekl.} & Helligkeit \\
\hline
April 12 & $335^\circ\;44'$ & $6^\circ\;52'$ & 0.0068 \\
16 & $336\;14$ & $7\;6$ & 0.0067 \\
20 & $336\;42$ & $7\;21$ & 0.0065 \\
24 & $337\;6$ & $7\;37$ & 0.0064 \\
28 & $337\;28$ & $7\;55$ & 0.0063 \\
Mai\ \ 2 & $337\;46$ & $8\;14$ & 0.0062 \\
6 & $338\;1$ & $8\;35$ & 0.0061 \\
10 & $338\;12$ & $8\;57$ & 0.0061 \\
14 & $338\;20$ & $9\;21\tfrac{1}{2}$ & 0.0060 \\
18 & $338\;23$ & $9\;48$ & 0.0060 \\
22 & $338\;23$ & $10\;16$ & 0.0059 \\
26 & $338\;18$ & $10\;46$ & 0.0059 \\
30 & $338\;10$ & $11\;19$ & 0.0058 \\
\end{tabular}
\end{center}
